\documentclass[12pt]{article}
\usepackage[utf8]{inputenc}
\usepackage[T1]{fontenc}
\usepackage{hyperref}
\title{Yogad Language Resources for K-12 Curriculum Development}
\author{Compiled by Claude Sonnet 5.5 (claude-sonnet-5-5) for Andr\'as Kornai}
\date{}
\begin{document}
\maketitle

\section{Introduction}
Yogad is an Austronesian language of the Philippines, spoken mainly in Echague and nearby towns in Isabela province. It is endangered, with roughly 16,000 speakers, and is not known to be taught in schools. The search found few materials. The main ones are a published dictionary and grammar and a 2021 orthography booklet from the Komisyon sa Wikang Filipino. I found no primers, textbooks or school programmes online.

\section{Primary Educational Resources}
No classroom textbooks or readers turned up in the search. Wikipedia mentions an SIL \emph{Yogad: First Primer} (1956), but I did not check it and give no link. The orthography booklet below is the closest thing to a teaching resource.
\begin{itemize}
\item \emph{Ortografiya nu Yogad} (Komisyon sa Wikang Filipino, 2021), the standard spelling guide, with a 21-letter alphabet: \url{https://kwfwikaatkultura.ph/wp-content/uploads/2024/03/Yogad.pdf}
\end{itemize}

\section{Linguistic and Reference Materials}
\begin{itemize}
\item Davis and Mesa, \emph{A Dictionary of Yogad} (LINCOM Europa, 2000). It is based on a speaker from Echague. The book is announced here, with no free full text: \url{https://linguistlist.org/issues/11/1367/}
\item Davis, Baker, Spitz and Baek, \emph{The Grammar of Yogad: A Functional Explanation} (1998). This is a printed book. No free online copy turned up.
\item Overview of the language: \url{https://en.wikipedia.org/wiki/Yogad_language}
\item Orthography comparison across Philippine languages: \url{https://en.wikipedia.org/wiki/Comparison_of_orthographies_of_languages_of_the_Philippines}
\end{itemize}

\section{Religious and Community Resources}
None found in the search. Wikipedia mentions a ``Yogad Kan'' website (yogadkan.com), but I did not see it in the results and did not verify it.

\section{Digital and Audio Resources}
\begin{itemize}
\item Short Yogad verb list with glosses (not a full dictionary): \url{https://docs.verbix.com/Languages/Yogad}
\item Online Yogad entry on a Tagalog dictionary site: \url{https://www.tagalog.com/dictionary/yogad}
\end{itemize}
No audio or video resources were found.

\section{Curriculum Development Context}
A curriculum would rest on the Davis and Mesa dictionary and grammar plus the KWF orthography. Both are scholarly or official reference works, not teaching materials. Primers, readers and recordings would have to be created with the community, and the available documentation is thin. Contact with the Komisyon sa Wikang Filipino and with Yogad speakers in Isabela would be the logical first step.
\end{document}